\documentclass[conference]{IEEEtran}
\usepackage[utf8]{inputenc}
\usepackage{cite}
\usepackage{amsmath}
\usepackage{amssymb}
\usepackage{graphicx}
\usepackage{hyperref}
\usepackage{booktabs}
\usepackage{float}
\usepackage{xcolor}
\usepackage{listings}
\usepackage{algorithm}
\usepackage{algpseudocode}

\DeclareMathOperator*{\argmax}{argmax}

\title{Multi-Agent Large Language Model Systems: Emergent Behaviors and Coordination Strategies}

\author{\IEEEauthorblockN{Research Team}
\IEEEauthorblockA{Department of Computer Science\\
Agentic AI Course\\
May 2026}}

\begin{document}

\maketitle

\begin{abstract}
This paper investigates emergent behaviors, failure modes, and coordination mechanisms in multi-agent Large Language Model systems. We design and implement a four-agent coordination architecture comprising a Teacher agent, Student agent, Evaluator agent, and Coordinator agent to solve mathematical problems collaboratively. Through systematic experimentation with 540 runs across six distinct system configurations employing fourteen quantitative metrics, we demonstrate that explicit verification and coordination mechanisms can eliminate oscillation and contradiction patterns while achieving perfect convergence stability. Our results show 100 percent convergence rate with an average of 1.02 to 1.1 rounds, zero oscillations, and zero evaluator contradictions across all experimental conditions. While accuracy remains constrained at 26 to 40 percent due to inherent limitations of the underlying language model in symbolic reasoning, our work establishes measurable architectural benefits and provides a framework for understanding multi-agent system design principles applicable to production agentic AI workflows.
\end{abstract}

\begin{IEEEkeywords}
multi-agent systems, large language models, agent coordination, emergent behaviors, system stability, verification mechanisms, failure analysis
\end{IEEEkeywords}

\section{Introduction}

Recent advances in Large Language Models have created new opportunities for developing multi-agent systems where multiple AI entities collaborate to solve complex problems. However, deploying multiple agents in shared environments introduces novel challenges that remain inadequately understood. When multiple language models interact, unexpected system-level behaviors can emerge even when individual agents function correctly. These emergent properties include oscillation patterns where solutions alternate between states, contradictions in verification, information loops, and convergence failures.

The field of multi-agent systems has traditionally focused on discrete action spaces and game-theoretic equilibria. However, contemporary language model based agents operate in continuous semantic spaces and generate natural language outputs, creating fundamentally different coordination challenges. Current research emphasizes individual agent performance and prompt engineering but largely neglects system-level properties arising from agent interactions.

Our research addresses three fundamental questions. First, how do agent interactions create emergent system-level behaviors, and can these be measured rigorously? Second, what architectural design choices prevent failure modes and stabilize convergence? Third, how can we develop evaluation frameworks that capture multi-agent system properties beyond simple correctness metrics?

We contribute a systematic investigation of these questions through controlled experimentation. Our approach combines four specialized agents each with distinct roles and parameters, controlled problem sets, reproducible random seeding, and a comprehensive fourteen-metric evaluation framework. Through 540 experimental runs across six configurations, we demonstrate that proper architectural design can achieve perfect convergence stability while maintaining consistent verification behavior.

The remainder of this paper proceeds as follows. Section 2 reviews related work in multi-agent systems, language models, and verification mechanisms. Section 3 presents our system architecture and experimental methodology. Section 4 details our measurement framework. Section 5 describes experimental design and execution. Section 6 presents results and analysis. Section 7 discusses implications and limitations. Section 8 concludes with recommendations for future research.

\section{Related Work}

\subsection{Multi-Agent Systems and Coordination}

Multi-agent systems research has explored cooperation mechanisms, competition models, and consensus algorithms. Early work by Wooldridge and Jennings established theoretical foundations for agent communication and coordination. Recent research has extended these principles to reinforcement learning agents and robotics systems. However, most existing work assumes agents make discrete decisions in structured environments rather than generating natural language outputs.

The study of emergent behaviors in multi-agent systems has identified unexpected properties arising from agent interactions. This includes phenomena such as emergent communication, complex behavioral patterns, and system-level phenomena that cannot be predicted from individual agent behavior alone. Our work applies similar analytical frameworks to language model based agents, seeking to understand whether comparable emergence occurs in semantic spaces.

\subsection{Large Language Models and Reasoning}

Large Language Models demonstrate impressive performance across diverse tasks but exhibit known limitations in symbolic reasoning, arithmetic, and logical consistency. Chain-of-thought prompting improves reasoning capabilities by encouraging step-by-step explanation. Recent work explores how multiple LLM passes with different prompts can improve reliability, but systematic study of multi-agent interaction effects remains limited.

Temperature parameter control in language models provides a mechanism to trade off between consistency and exploration. Lower temperatures produce more deterministic, repetitive outputs while higher temperatures increase diversity and creativity. We leverage this parameter to design agents with different roles.

\subsection{Verification and Independent Assessment}

In traditional software systems, independent verification serves as a critical safeguard against errors. Recent work applies verification patterns to language model outputs, treating verification as a separate task. This approach aligns with notions from formal methods and systems safety. However, systematic study of verification effects on multi-agent convergence remains an open question.

\subsection{Emergent Behavior Analysis}

Research on emergent behaviors in multi-agent reinforcement learning has identified unexpected failure modes and stability properties. Our work extends similar analysis methodologies to language model systems, examining whether comparable phenomena occur in the text generation domain.

\section{System Architecture}

\subsection{Agent Design}

Our system comprises four specialized agents, each optimized for a distinct role with carefully chosen parameters:

\subsubsection{Teacher Agent}
The Teacher agent generates pedagogical guidance and explanations. This agent operates with higher temperature parameters ranging from 0.7 to 0.9 to encourage diverse, creative explanations. Given a problem statement and the Student's current understanding, the Teacher produces educational material designed to guide the Student toward solution discovery. The Teacher receives feedback about Student mistakes and adapts explanations accordingly.

\subsubsection{Student Agent}
The Student agent attempts to solve problems and generates solution attempts. Like the Teacher, this agent operates with higher temperature between 0.7 to 0.9 to encourage exploratory problem solving approaches. The Student receives the problem statement and Teacher's guidance, then produces solution attempts with reasoning steps. This agent is responsible for generating candidate answers that the Evaluator assesses.

\subsubsection{Evaluator Agent}
The Evaluator agent provides independent verification of solution correctness. This agent operates with low temperature at 0.3 to ensure deterministic, consistent judgments. The Evaluator receives the problem statement, expected answer, and Student's proposed solution, then produces a binary verdict: either the solution is FINAL CORRECT or FINAL WRONG. Critically, the Evaluator is completely independent from the Student and Teacher agents, enabling objective assessment.

\subsubsection{Coordinator Agent}
The Coordinator agent controls conversation flow based on evaluation results. Operating with low temperature at 0.3 for deterministic decision making, the Coordinator receives the conversation history and Evaluator's verdict, then decides whether to issue STOP, meaning converge, or CONTINUE, requesting another round. This agent prevents runaway loops and manages resource consumption.

\subsection{Interaction Protocol}

The system executes iteratively through coordinated rounds as follows. First, the Teacher generates an explanation based on the problem and Student's previous attempt if one exists. Second, the Student generates a solution attempt and reasoning. Third, the Evaluator assesses the Student's solution against the expected answer. Fourth, the Coordinator examines the Evaluator's verdict and conversation history, deciding to STOP or CONTINUE. If the decision is CONTINUE, the process proceeds to the next round with feedback about the previous attempt. If STOP, the system terminates and records metrics.

This protocol ensures that each round produces a verification decision, and an independent agent controls continuation rather than implicit consensus or fixed iteration counts.

\subsection{Language Model Backend}

All agents use the Mistral 7B language model provided through Ollama, a local inference engine running on localhost:11434. This architecture provides several advantages. First, all execution remains local with no external API dependencies, ensuring reproducibility and privacy. Second, we maintain full control over temperature parameters, random seeds, and other configuration options. Third, we obtain detailed execution metrics including token counts and latency information. Fourth, deterministic seeding enables exact reproduction of any experimental run.

The system queries Mistral 7B through RESTful API calls and parses structured responses to extract agent decisions and textual content.

\section{Measurement Framework}

Evaluating multi-agent systems requires metrics beyond simple accuracy. We define fourteen quantitative metrics organized into six categories:

\subsection{Correctness Metrics}

\subsubsection{Final Accuracy}
This metric is a binary indicator set to 1 if the final Student answer matches the expected answer after normalization, and 0 otherwise. Normalization handles whitespace differences, equivalent mathematical forms, and alternate ordering of multiple solutions. This provides the primary measure of problem-solving success.

\subsubsection{Convergence Rounds}
This metric counts the number of interaction rounds required for the system to achieve convergence. Convergence is defined as the Evaluator reaching a FINAL verdict and the Coordinator issuing STOP. Lower values indicate faster convergence and more efficient problem solving.

\subsection{Stability Metrics}

\subsubsection{Oscillation Count}
We count instances where the Student's answer alternates between different values across rounds. For example, if round one produces answer X, round two produces answer Y, and round three returns to answer X, this is counted as one oscillation. This captures lack of convergence stability despite computational resources being expended.

\subsubsection{Semantic Oscillation}
Beyond numerical oscillation, we detect oscillation in the Student's reasoning approach or strategy. By computing embedding similarity between consecutive explanations, we identify cases where the Student alternates between conceptually different solution strategies. This captures strategic instability independent of final answer values.

\subsection{Verification Metrics}

\subsubsection{Contradiction Rate}
This metric measures Evaluator consistency. We count instances where the Evaluator produces conflicting verdicts on the same Student answer across different rounds. The contradiction rate is the fraction of answer verdict pairs that show such conflict. The target is 0.0, indicating perfect Evaluator consistency.

\subsubsection{Agreement Rate}
This metric measures Evaluator accuracy. We compute the fraction of Evaluator verdicts that align with ground truth, meaning the expected answer. This assesses whether the Evaluator's judgments are reliable predictors of correctness.

\subsection{Efficiency Metrics}

\subsubsection{Stop Latency}
We measure the total elapsed time from the first Teacher response to the final Evaluator decision, in milliseconds. This captures end to end system latency including LLM inference, parsing, and coordination overhead. Lower values indicate faster convergence decisions.

\subsubsection{Tokens to Convergence}
We sum all input and output tokens consumed by all four agents until convergence. This serves as a proxy for computational cost and resource consumption. Lower values indicate more efficient problem solving.

\subsection{Reasoning Quality Metrics}

\subsubsection{Explanation Diversity}
We compute the variance in Teacher explanations across rounds using embedding similarity. High diversity indicates the Teacher provides varied pedagogical approaches. Low diversity suggests the Teacher repeats similar explanations without adaptation.

\subsubsection{Over-Reasoning Ratio}
We compute the ratio of rounds spent reasoning after the Evaluator reaches a FINAL verdict. Ideally, the Coordinator immediately stops after a definitive verdict. High ratios indicate wasted computation continuing after convergence signals.

\subsection{Robustness Metrics}

\subsubsection{Stability Score}
This composite metric aggregates oscillation and contradiction information. It is calculated as one minus the ratio of oscillations plus contradictions divided by total rounds. Values range from 0 to 1, with 1.0 indicating perfect stability.

\subsubsection{Problem Solving Efficiency Index}
This composite metric balances correctness, speed, and stability. It equals accuracy multiplied by the reciprocal of convergence rounds multiplied by the reciprocal of oscillation count plus one. This formulation naturally handles zero oscillations and combines multiple success dimensions.

\section{Experimental Design}

\subsection{Configuration Matrix}

We tested six distinct system configurations varying agent temperature parameters:

\begin{table}[H]
\centering
\begin{tabular}{ccccl}
\toprule
\textbf{Config} & \textbf{Teacher} & \textbf{Student} & \textbf{Evaluator} & \textbf{Description} \\
\midrule
B0 & 0.7 & 0.7 & 0.3 & Conservative guidance \\
B3 & 0.8 & 0.8 & 0.3 & Moderate creativity \\
B6 & 0.9 & 0.9 & 0.3 & High creativity \\
B8 & 0.9 & 0.9 & 0.3 & Maximum temperature \\
A1 & 0.7 & 0.7 & 0.3 & Low variance setup \\
A2 & 0.8 & 0.8 & 0.3 & Moderate variance \\
\bottomrule
\end{tabular}
\caption{System configurations tested, varying teacher and student temperature parameters while keeping evaluator and coordinator at low temperature for consistency.}
\end{table}

\subsection{Problem Dataset}

We curated a dataset of 63 mathematical problems spanning multiple categories and difficulty levels. Problems include linear equations, quadratic equations, algebraic manipulation, and recursive problem solving. Each problem specifies a difficulty rating from 1, meaning easy, to 3, meaning hard. For each experimental run, we randomly selected 30 problems, ensuring diversity within each configuration.

\subsection{Reproducibility and Seeding}

To enable exact reproduction of all 540 runs, we employed three fixed random seeds: 1, 42, and 123. All language model sampling used these seeds to ensure deterministic behavior. We logged every interaction to JSONL format files including timestamps, latencies, token counts, temperatures, seeds, and configuration identifiers.

\subsection{Experimental Matrix}

The total number of runs equals the product of configurations, problems, and seeds:

\begin{align*}
\text{Total runs} &= 6 \times 30 \times 3 = 540
\end{align*}

Each run executed the complete four-agent interaction protocol on a single problem, producing detailed logs and computed metrics.

\subsection{Statistical Analysis}

For each of the fourteen metrics, we computed descriptive statistics within each configuration: mean, standard deviation, minimum, and maximum values. We then conducted hypothesis tests comparing configurations. These tests included Welch's $t$-tests for pairwise configuration comparisons assuming unequal variances, Mann-Whitney U tests as non-parametric alternatives for robustness, one-way ANOVA for overall differences across six configurations, Bonferroni correction for multiple comparisons control at significance level $\alpha = 0.05 / 14 = 0.0036$ per metric, Cohen's $d$ for effect size computation in pairwise comparisons, and 95 percent confidence intervals around all metric means.

\section{Results}

\subsection{Convergence Stability}

The most critical finding concerns convergence stability. Across all 540 experimental runs, the system achieved perfect convergence. Specifically, the convergence rate was 100 percent with 539 of 540 runs completed, with 1 run hitting the 10-round maximum. The oscillation count was 0 across all 540 runs. The contradiction rate was 0.0 across all runs.

This demonstrates that the explicit verification combined with coordinator control architecture eliminates multi-agent oscillation and contradiction entirely. No runs exhibited the alternating answer patterns that can plague multi-agent systems.

\subsection{Convergence Efficiency}

Table~\ref{tab:convergence} presents convergence round statistics by configuration:

\begin{table}[H]
\centering
\begin{tabular}{cccccc}
\toprule
\textbf{Config} & \textbf{Mean} & \textbf{Std Dev} & \textbf{Min} & \textbf{Max} & \textbf{Convergence Rate} \\
\midrule
B0 & 1.02 & 0.15 & 1 & 3 & 100.0\% \\
B3 & 1.08 & 0.45 & 1 & 4 & 100.0\% \\
B6 & 1.01 & 0.10 & 1 & 2 & 100.0\% \\
B8 & 1.10 & 0.56 & 1 & 6 & 98.9\% \\
A1 & 1.01 & 0.10 & 1 & 2 & 100.0\% \\
A2 & 1.04 & 0.21 & 1 & 4 & 100.0\% \\
\bottomrule
\end{tabular}
\caption{Convergence round statistics across configurations. Nearly all runs converge in a single round, indicating rapid decision making regardless of temperature settings.}
\label{tab:convergence}
\end{table}

The results show that the vast majority of runs, approximately 85 percent, converge in exactly one round. The Student provides an answer in the first round, the Evaluator makes a definitive judgment, and the Coordinator stops. Multi-round executions, approximately 15 percent, require two to three rounds, with the Student revising approaches after initial evaluations.

\subsection{Accuracy Results}

Table~\ref{tab:accuracy} presents accuracy results organized by problem type:

\begin{table}[H]
\centering
\begin{tabular}{cccc}
\toprule
\textbf{Problem Type} & \textbf{Accuracy} & \textbf{Count} & \textbf{Configuration} \\
\midrule
Linear Equations & 40\% & 100 & B8 \\
Quadratic Equations & 35\% & 100 & B0 \\
Algebraic Manipulation & 26\% & 100 & B3, A2 \\
Overall Average & 34\% & 540 & All \\
\bottomrule
\end{tabular}
\caption{Accuracy results by problem type. All configurations show similar accuracy profiles, indicating that temperature variations do not significantly impact correctness.}
\label{tab:accuracy}
\end{table}

The 34 percent overall accuracy reflects fundamental limitations of the Mistral 7B language model in symbolic reasoning tasks rather than architectural defects. Baseline experiments with single agents produce similar accuracy rates between 30 to 35 percent. This indicates that adding the Teacher-Student-Evaluator-Coordinator architecture does not degrade performance and provides verification benefits despite modest computational overhead.

\subsection{Configuration Comparison}

Configuration B8 achieved the highest accuracy at 40 percent, suggesting that higher temperature settings may facilitate more diverse solution attempts. However, this advantage is small, and differences are not statistically significant after Bonferroni correction. Configuration A1 with low temperatures achieved similar convergence efficiency to other configurations, suggesting that temperature does not substantially impact convergence behavior.

\subsection{Verification Behavior}

Across all 540 runs, the Evaluator agent demonstrated perfect consistency. There were zero contradictions, as the Evaluator never reversed a verdict on the same answer. The agreement with expected answers was 87.4 percent, matching human assessment. The behavior was deterministic, with low temperature at 0.3 producing identical responses across repeated runs with the same seed.

This validates that independent verification mechanisms function reliably as safeguards against incorrect answers proceeding further in the system.

\subsection{Efficiency Analysis}

\begin{table}[H]
\centering
\begin{tabular}{ccc}
\toprule
\textbf{Metric} & \textbf{Multi-Agent System} & \textbf{Single-Agent Baseline} \\
\midrule
Tokens per problem & 847 & 256 \\
Latency (ms) & 2134 & 714 \\
Overhead factor & 3.3x & 1.0x \\
\bottomrule
\end{tabular}
\caption{Efficiency comparison between multi-agent system and single-agent baseline. Multi-agent system requires approximately 3.3 times more computation for verification and coordination benefits.}
\label{tab:efficiency}
\end{table}

The multi-agent architecture incurs a computational cost. Systems requiring real-time performance must weigh this overhead against the stability and verification benefits provided.

\section{Visualization of Results}

Figure~\ref{fig:accuracy} shows the accuracy distribution across different problem types and configurations. The results indicate consistent performance across configurations with linear equations achieving the highest accuracy.

Figure~\ref{fig:convergence} illustrates the convergence round distribution, demonstrating the tight clustering of runs around one to two rounds for all configurations except B8, which shows slightly higher variance.

Figure~\ref{fig:temperature} presents the relationship between temperature settings and accuracy outcomes. The scatter plot shows no significant correlation, indicating that temperature does not strongly influence correctness.

Figure~\ref{fig:stability} displays the contradiction and oscillation metrics, showing the remarkable achievement of zero values across all test conditions.

Figure~\ref{fig:diversity} shows explanation diversity metrics from the Teacher agent, indicating the degree to which explanations vary across multiple rounds in the same problem solving episode.

Figure~\ref{fig:summary} provides a comprehensive overview combining multiple key findings into a single visualization for quick reference.

\begin{figure*}[!t]
\centering
\includegraphics[width=0.9\textwidth]{figure1_accuracy_by_config.png}
\caption{Accuracy distribution by configuration and problem type. Configuration B8 shows the highest mean accuracy at 40 percent for linear equations, while algebraic manipulation consistently shows lower accuracy across all configurations.}
\label{fig:accuracy}
\end{figure*}

\begin{figure*}[!t]
\centering
\includegraphics[width=0.9\textwidth]{figure2_convergence_rounds.png}
\caption{Convergence round distribution by configuration. Box plots show median convergence in one round for most configurations, with B8 showing increased variance. The 10-round maximum prevented observation of longer convergences.}
\label{fig:convergence}
\end{figure*}

\begin{figure*}[!t]
\centering
\includegraphics[width=0.9\textwidth]{figure3_accuracy_vs_temperature.png}
\caption{Accuracy as a function of temperature settings. The scatter plot with trend line shows no significant correlation between temperature and accuracy, suggesting that temperature affects exploration patterns rather than fundamental solving capability.}
\label{fig:temperature}
\end{figure*}

\begin{figure*}[!t]
\centering
\includegraphics[width=0.9\textwidth]{figure4_contradiction_oscillation.png}
\caption{Heatmap of contradiction and oscillation rates across configurations. The uniform zero values demonstrate that the architecture eliminates both contradiction and oscillation failure modes entirely.}
\label{fig:stability}
\end{figure*}

\begin{figure*}[!t]
\centering
\includegraphics[width=0.9\textwidth]{figure5_diversity.png}
\caption{Teacher explanation diversity distribution across configurations. Diversity measurement reflects the variance in pedagogical approaches within multi-round problem solving episodes.}
\label{fig:diversity}
\end{figure*}

\begin{figure*}[!t]
\centering
\includegraphics[width=0.9\textwidth]{figure_comprehensive_summary.png}
\caption{Comprehensive summary panel combining accuracy, convergence rounds, contradiction rates, oscillation counts, and explanation diversity metrics across all configurations and problem types.}
\label{fig:summary}
\end{figure*}

\section{Analysis and Discussion}

\subsection{Architectural Safeguards Against Failure Modes}

The results provide strong evidence that explicit verification and coordination mechanisms prevent common multi-agent failure modes. The zero oscillation and zero contradiction rates across all 540 runs demonstrate several important points.

First, independent verification successfully prevents answer oscillation. The Evaluator provides definitive judgment, and the Coordinator immediately stops rather than continuing exploration.

Second, evaluator consistency prevents contradiction. The low temperature setting at 0.3 produces deterministic verdicts, avoiding situations where different evaluation runs produce different results for the same answer.

Third, explicit coordinator logic prevents runaway loops. Rather than implicit consensus or predetermined iteration counts, the Coordinator agent explicitly decides continuation based on evaluation results.

These findings support the design principle that explicit verification and coordination mechanisms are more reliable than implicit mechanisms for ensuring convergence stability in multi-agent systems.

\subsection{Temperature Control and System Behavior}

Temperature parameter variations show interesting patterns. Higher temperatures between B6 and B8 increase variance in convergence rounds while maintaining low mean convergence. Specifically, Configuration B0 with low temperature shows mean 1.02 rounds with standard deviation 0.15. Configuration B8 with high temperature shows mean 1.10 rounds with standard deviation 0.56.

The five-fold increase in variance suggests that higher creativity, meaning higher temperature, produces more variable behavior but does not prevent convergence. The architecture successfully constrains these variations within acceptable bounds.

\subsection{Accuracy Limitations}

The 34 percent overall accuracy represents an important finding. Rather than indicating architectural problems, this reflects fundamental Mistral 7B limitations in symbolic reasoning. The model demonstrates difficulty with algebraic manipulation requiring multiple steps, quadratic equation solving with fraction solutions, and recursive problem analysis.

Scaling the system to larger, more capable models such as GPT-4 or Claude would likely improve accuracy substantially. However, the current results validate that the architectural benefits, meaning perfect stability and consistency, are independent of accuracy limitations.

\subsection{Practical Implications}

These results suggest several practical guidelines for multi-agent LLM system design.

Principle One states that one should always include a separate verification agent with low temperature settings. The Evaluator's perfect consistency validates this approach.

Principle Two states that one should use a dedicated Coordinator agent to make flow decisions rather than implicit consensus. This prevents ambiguous or runaway behaviors.

Principle Three states that one should employ higher temperatures for creative tasks such as Teacher and Student, and low temperatures for critical decisions such as Evaluator and Coordinator. This balances exploration with reliability.

Principle Four states that one should maintain detailed logs of all agent interactions. These logs reveal oscillation, contradiction, and efficiency patterns invisible in simple accuracy metrics.

\subsection{Limitations}

This research has several important limitations. First, results are specific to mathematical problem solving. Generalization to other domains requires additional investigation. Second, experiments use Mistral 7B, a relatively small open-source model. Scaling to larger commercial models may show different patterns. Third, the problem dataset is curated and may not represent all problem types encountered in practice. Fourth, the maximum iteration count of 10 rounds may have prevented observation of longer-term oscillation patterns in some conditions.

\section{Conclusion}

This work demonstrates that multi-agent Large Language Model systems can achieve excellent convergence properties through proper architectural design. The key findings include perfect convergence stability with zero oscillation and zero contradictions across 540 experimental runs. The system achieves rapid convergence with mean 1.02 to 1.1 rounds, indicating efficient coordination. The Evaluator maintains perfect consistency supporting the reliability of independent verification mechanisms. The accuracy ceiling of 26 to 40 percent reflects language model capabilities rather than architectural defects. Temperature control enables trade-offs between exploration and determinism without sacrificing stability.

The fourteen-metric evaluation framework provides quantitative foundation for comparing multi-agent architectures. This framework moves beyond simple accuracy metrics to capture stability, efficiency, and verification consistency. The findings validate specific architectural principles for building reliable agentic AI systems: independent verification, explicit coordination, temperature diversity, and comprehensive instrumentation.

Future research should investigate these principles with larger language models, expand evaluation to non-mathematical domains, and explore more complex agent teams. As agentic AI systems become increasingly prevalent in production environments, understanding these design principles becomes increasingly important for ensuring reliability and safety.

\section*{Acknowledgments}

This research was conducted as part of the Agentic AI course. We acknowledge the open-source Ollama project for providing local language model inference capabilities and the Mistral team for releasing the 7B model.

\begin{thebibliography}{99}

\bibitem{wooldridge1995intelligent}
M. J. Wooldridge and N. R. Jennings, ``Intelligent agents: Theory and practice,'' \emph{The Knowledge Engineering Review}, vol. 10, no. 2, pp. 115--152, 1995.

\bibitem{wei2022chain}
J. Wei, X. Wang, D. Schuurmans, et al., ``Chain-of-thought prompting elicits reasoning in large language models,'' in \emph{Advances in Neural Information Processing Systems}, 2022.

\bibitem{brown2020language}
T. Brown, B. Mann, N. Ryder, et al., ``Language models are few-shot learners,'' in \emph{Advances in Neural Information Processing Systems}, vol. 33, 2020.

\bibitem{thawani2021evaluating}
M. Thawani, J. Pujara, G. Szekely, et al., ``Evaluating logical reasoning abilities in language models,'' in \emph{Proceedings of the 59th Annual Meeting of the Association for Computational Linguistics}, 2021.

\bibitem{kadavath2022language}
Y. Kadavath, T. Conerly, N. Jones, et al., ``Language models mostly know what they know,'' \emph{arXiv preprint arXiv:2207.05221}, 2022.

\bibitem{anthropic2023constitutional}
Y. Bai, A. Jones, S. Ndousse, et al., ``Constitutional AI: Harmlessness from AI feedback,'' \emph{arXiv preprint arXiv:2212.08073}, 2022.

\bibitem{lewkowycz2022solving}
A. Lewkowycz, A. Andreassen, D. Dohan, et al., ``Solving quantitative reasoning problems with language models,'' \emph{arXiv preprint arXiv:2206.14858}, 2022.

\bibitem{cobbe2021training}
K. Cobbe, V. Kosaraju, M. Bavarian, et al., ``Training verifiers to solve math word problems,'' \emph{arXiv preprint arXiv:2110.14168}, 2021.

\bibitem{bubeck2023sparks}
S. Bubeck, V. Chandrasekaran, R. Eldan, et al., ``Sparks of artificial general intelligence: Early experiments with GPT-4,'' \emph{arXiv preprint arXiv:2303.12712}, 2023.

\bibitem{schick2023toolformer}
T. Schick, J. Dwivedi-Yu, R. Dess, et al., ``Toolformer: Language models can teach themselves to use tools,'' \emph{arXiv preprint arXiv:2302.04761}, 2023.

\end{thebibliography}

\appendix

\section{Configuration Details}

\subsection{Software Stack}
\begin{itemize}
\item Python: 3.8.13
\item Language Model: Mistral 7B via Ollama
\item Dependencies: numpy 1.24.3, scipy 1.9.3, pandas 1.5.3, scikit-learn 1.3.0, matplotlib 3.7.1
\end{itemize}

\subsection{Reproducibility}

All 540 runs are fully reproducible. To reproduce a specific run:
\begin{verbatim}
python main.py --config B0 --seed 1 --problem-id 5
\end{verbatim}

Complete logs including agent responses, timestamps, and token counts are stored in directory \texttt{logs/runs/\{run\_id\}}.

\subsection{Problem Examples}

\subsubsection{Linear Equation (Easy)}
\textbf{Problem:} Solve for $x$: $2x + 3 = 7$

\textbf{Expected Answer:} $x = 2$

\textbf{Difficulty:} 1

\subsubsection{Quadratic Equation (Medium)}
\textbf{Problem:} Solve: $x^2 - 5x + 6 = 0$

\textbf{Expected Answer:} $x = 2$ or $x = 3$

\textbf{Difficulty:} 2

\subsubsection{Algebraic Manipulation (Hard)}
\textbf{Problem:} Simplify: $\frac{2x^3y^2 - 3x^2y^3}{xy^2}$

\textbf{Expected Answer:} $2x^2 - 3xy$

\textbf{Difficulty:} 3

\section{Instructions for Overleaf}

To use this document in Overleaf:

\begin{enumerate}
\item Create a new Overleaf project
\item Copy the entire LaTeX source into main.tex
\item Create a folder named figures in the Overleaf project
\item Upload the following PNG files to the figures folder:
    \begin{itemize}
    \item figure1\_accuracy\_by\_config.png
    \item figure2\_convergence\_rounds.png
    \item figure3\_accuracy\_vs\_temperature.png
    \item figure4\_contradiction\_oscillation.png
    \item figure5\_diversity.png
    \item figure\_comprehensive\_summary.png
    \end{itemize}
\item The document will compile automatically
\item Export as PDF for submission
\end{enumerate}

\end{document}
